\section{Task inventory and implementation details}
\label{app:tasks}
PU uses BCD$\to$A, ACD$\to$B, ABD$\to$C and ABC$\to$D. JNU uses
BC$\to$A, AC$\to$B and AB$\to$C. PHM09 adjacent-domain tasks are
AC$\to$B, BC$\to$A, BD$\to$C, CE$\to$D, DF$\to$E,
EG$\to$F, FH$\to$G, GI$\to$H, HI$\to$J and HJ$\to$I.
Its distant-domain tasks are AB$\to$F/G/H/I/J and IJ$\to$A/B/C/D/E,
where each target letter defines a separate task.

Source context length is 32768 samples and response support is 4 ms.
The four carrier candidates per context are selected from source spectra above
four times the largest paired source shaft rate and below 0.44 times the
sampling frequency. Candidate peaks are separated by at least
$\max(100,3f_s/N)$ Hz. Each carrier is combined with decay constants
0.25, 0.5 and 1 ms and phases 0, $\pi/2$, $\pi$ and $3\pi/2$.
The kernel is an exponentially decaying sinusoid multiplied by a squared
cosine taper and normalized by its discrete $\ell_2$ norm.
Numerical coefficients are obtained by the fixed accelerated proximal solver;
the classifier does not update them.
The reused numerical fits use 500 proximal iterations. Their penalties are
specified from each normalized source context. Let
\begin{equation}
 \nu_j=\max\!\left\{
 \frac{\operatorname{median}_{f\in[0.45,0.5]}
       |\operatorname{rFFT}(\bar x_j)[f]|}{\sqrt{N\log 2}},
 0.005\sqrt{\frac1N\sum_t\bar x_j[t]^2}\right\},
\end{equation}
where $f$ is normalized frequency, and define
\begin{equation}
 b(G)=\frac{\nu_1+\nu_2}{2}
 \sqrt{4+4\sqrt{\log(\max(G,2))}+2\log(\max(G,2))}.
\end{equation}
Then $\lambda_{s,j}=\nu_j\sqrt{2\log(2\cdot48N)}$,
$\lambda_b=b(4)$ and $\lambda_c=b(\max(G_F,1))$, where $G_F$ is the number
of stationary frequency candidates. The proximal step uses the fixed
spectral-norm upper bound formed from the convolution, shaft-harmonic and
stationary dictionaries. These numerical settings are source-derived and
are reused for all edited views of a context.
